\section{Deployment lessons}\label{sec:deployment_lessons}
%We started deployment of {\systemname} with offline inference and observed that we are loosing on some of the fresh candidates, in particular we in A/B test we noticed that live updates plays a key role to serve freshly created posts at LinkedIn. We observed that live-updates can bring relative gains of +6\% in our production systems when live-update is enabled. That is why live updates are so important for improved system performance.

\textit{Freshness:} We initially deployed {\systemname} with offline inference and found it missed some fresh candidates. A/B tests revealed that live updates are crucial for serving newly created LinkedIn posts. Enabling live updates resulted in a +6\% gain in our production systems, highlighting their importance for improved performance.

\textit{Pre-filtering}: Our system employs EBR with pre-filtering, significantly enhancing retrieval quality. Many existing EBR infrastructures use KNN search with post-filtering, where results are first retrieved by KNN distance and then filtered. This post-filtering approach reduces system recall and quality by wasting candidate slots on items that don't meet attribute constraints. By enabling pre-filtering on GPU retrieval, we greatly improved the quality of results compared to our production FAISS and lucene-based systems.

\textit{Custom filtering kernel}: One lesson we learned early is that native TF or PyTorch do not effectively support filtering operations because deep learning frameworks weren't initially designed for model-based retrieval indexes. Native boolean masking and indexing cause a 100X latency increase, making them impractical for production. Therefore, we implemented a custom CUDA solution for pre-filtering on the GPU, which scans items in memory to find those that meet constraints. One approach is to create a CUDA filtering kernel and fuse it with the matrix multiplication kernel to perform masked-matrix multiplication for KNN with pre-filtering. However, this solution is hard to generalize to other similarity measures or operations, as each new architecture would require re-implementation and fine-tuning, slowing down model development and deployment. 
In practice, we could make a trade-off of fully fused kernels and separately implemented kernels. For products needing regular KNN support, fusing the entire kernel with top-k selection and quantization improves serving speed. For general use cases, we create individual custom operations, like pre-filtering and quantization, to allow flexible development and deployment of advanced selection strategy with native neural network operations supported by TF and PyTorch. It is noteworthy that all the results reported in the paper are based on the second solution (even for the three plain KNN version) without extra custom kernel fusion for the purpose of self-consistency and generalizability.
%One of the lessons we learnt early on that native TF or PyTorch does not effectively support filtering operations, as developers of deep learning frameworks didn't think initially about model-based retrieval indexes. Native boolean masking and indexing filtering will cause 100X latency, which is not feasible to be used in production applications. Therefore, we had to implement custom CUDA based solution to support pre-filtering of results on GPU, which can scan over items in GPU memory to effectively find subset of items that satisfy constraints. One feasible way of implementing this is to create a CUDA filtering kernel and fuse it with matrix multiplication kernel to create a masked-matrix multiplication kernel for regular KNN multiplication with pre-filtering. However, this solution could be hard to generalize to universal similarity measurement such as a network-based similarity or plain hadamard product since we need to re-implement the corresponding kernel and fine-tune it every time a new architecture is created, which largely reduces the ML engineers' model development and deployment speed and lose the benefit of leveraging existing operations and kernel fusion ability of TF and PyTorch. In practice, we could make a trade-off of fully fused kernels and separately implemented kernels. For products that only needs regular KNN support, we could fuse the entire kernel (as well as the top-k selection and quantization operations) as one to potentially maximize the serving speed, and for general use cases, we create the individual custom operations, such as pre-filtering operation and quantization operation, to allow flexibility on developing and deploying advanced selection strategy with native neural network operations supported by TF and PyTorch. It is noteworthy that all the results reported in the paper are based on the second solution (even for the three plain KNN version) without extra custom kernel fusion for the purpose of self-consistency and generalizability.